\chapter{Topic: How Virtual Methods Work (Vtables and Object Layout)}

\begin{aside}
Material from the remainder of the term, not covered in the original lecture
notes.
\end{aside}

\section{The Question}
\begin{lstlisting}[style=cpp]
class Vec {
    int x, y;
    void f();
};

class Vec2 {
    int x, y;
    virtual void f();
};

Vec v{1, 2};
Vec2 w{1, 2};
cout << sizeof(v) << " " << sizeof(w);   // ?
\end{lstlisting}
On a typical 64-bit machine this prints \code{8 16}. Both have two \code{int}s
(8 bytes); methods are not stored in objects at all (there is one copy of each
method's code). So why is \code{w} bigger? Because of the virtual method.

\section{Dynamic Dispatch Needs Run-Time Information}
\begin{lstlisting}[style=cpp]
Book *pb = (userSaysComic ? new Comic{...} : new Book{...});
pb->isHeavy();   // Book::isHeavy or Comic::isHeavy?
\end{lstlisting}
The compiler can't know which method to call; the decision must be made at run
time, so the \emph{object} has to carry enough information to find the right
method.

\section{Vtables and Vptrs}
For each class with virtual methods, the compiler creates one \textbf{virtual
table} (\textbf{vtable}): a static array of pointers to the class's virtual
methods (plus the RTTI used by \code{dynamic\_cast}).

Each object of such a class contains a hidden pointer, the \textbf{vptr}, to its
class's vtable. The vptr is set by the constructor.

\begin{lstlisting}[style=cpp]
class Book {
    std::string title, author;
    int numPages;
  public:
    virtual bool isHeavy() const;
    virtual ~Book();
};
class Comic : public Book {
    std::string hero;
  public:
    bool isHeavy() const override;
};
\end{lstlisting}

\begin{center}
\begin{tikzpicture}[f/.style={draw, minimum width=25mm, minimum height=6mm,
  font=\small}, >={Stealth}]
% Book object
\node[font=\bfseries\small] at (0,0.9) {Book object};
\node[f] (bv) at (0,0.3) {vptr};
\node[f, below=0mm of bv] (bt) {title};
\node[f, below=0mm of bt] (ba) {author};
\node[f, below=0mm of ba] {numPages};
% Comic object
\node[font=\bfseries\small] at (0,-3.2) {Comic object};
\node[f] (cv) at (0,-3.8) {vptr};
\node[f, below=0mm of cv] (ct) {title};
\node[f, below=0mm of ct] (ca) {author};
\node[f, below=0mm of ca] (cn) {numPages};
\node[f, below=0mm of cn] {hero};
% vtables
\node[font=\bfseries\small] at (6,0.9) {Book vtable};
\node[f, minimum width=34mm] (bvt1) at (6,0.3) {RTTI: ``Book''};
\node[f, minimum width=34mm, below=0mm of bvt1] (bvt2) {isHeavy};
\node[f, minimum width=34mm, below=0mm of bvt2] {\textasciitilde Book};
\node[font=\bfseries\small] at (6,-3.2) {Comic vtable};
\node[f, minimum width=34mm] (cvt1) at (6,-3.8) {RTTI: ``Comic''};
\node[f, minimum width=34mm, below=0mm of cvt1] (cvt2) {isHeavy};
\node[f, minimum width=34mm, below=0mm of cvt2] {\textasciitilde Comic};
\draw[->] (bv.east) -- (bvt1.west);
\draw[->] (cv.east) -- (cvt1.west);
\node[right=3mm of bvt2, font=\scriptsize] {$\to$ \code{Book::isHeavy}};
\node[right=3mm of cvt2, font=\scriptsize] {$\to$ \code{Comic::isHeavy}};
\end{tikzpicture}
\end{center}

Calling a virtual method \code{pb->isHeavy()}:
\begin{enumerate}
  \item follow \code{pb} to the object;
  \item follow the object's vptr to its class's vtable;
  \item fetch the function pointer stored in the \code{isHeavy} slot (whose
        index is known at compile time);
  \item call it.
\end{enumerate}
This happens at run time, and it selects the right method no matter what the
declared type of the pointer is.

\section{Costs}
\begin{itemize}
  \item \textbf{Space}: one vptr (8 bytes) per object, plus one vtable per class.
  \item \textbf{Time}: two extra pointer dereferences per virtual call, and
        virtual calls usually can't be inlined.
\end{itemize}
That is why C++ makes methods non-virtual by default: you don't pay for dynamic
dispatch unless you ask for it. It is also why a class with no virtual methods
has exactly the same layout as a C struct.

\section{Object Layout and Inheritance}
Why does \code{Book *pb = \&someComic;} work at all? The compiler lays out a
subclass object as \emph{the superclass part first}, followed by the subclass's
own fields (the vptr, if any, comes first of all, as in g++). So a pointer to a
Comic is also a valid pointer to its Book part: the Book fields are at the same
offsets in both, and code compiled for Book can access them.

This also explains slicing: \code{Book b = comic;} copies only the first part of
the Comic (the part that fits in a Book), and \code{b}'s vptr is set by Book's
constructor to Book's vtable --- so \code{b.isHeavy()} calls
\code{Book::isHeavy}.

\section{Calling Virtual Methods in Constructors}
During Book's constructor (step 2 of constructing a Comic), the vptr points to
\emph{Book's} vtable; only once Comic's constructor starts is it switched to
Comic's. So a virtual call made inside a base class constructor (or destructor)
does \emph{not} dispatch to the subclass override: the subclass part isn't built
yet (or has already been destroyed).
\begin{important}[Rule]
Don't rely on virtual dispatch inside constructors and destructors.
\end{important}

\section*{Summary}
\begin{itemize}
  \item Each class with virtual methods has a vtable of method pointers; each
        object of such a class has a vptr to its class's vtable.
  \item A virtual call follows object $\to$ vptr $\to$ vtable slot $\to$ code.
  \item Cost: space for a vptr per object and an indirection per call; hence
        non-virtual is the default.
  \item Subclass objects start with the superclass part, which makes upcasting
        free.
  \item In ctors/dtors, virtual calls use the class currently being
        constructed/destroyed.
\end{itemize}
